\subsection{Actual poles at roots of exponential order in ordinary plethysm}

\begin{theorem}[P-F31: the chain pole refutes the positive-form clause of PH3]
\label{thm:actualcubicpole}
Let $n\ge4$, put
\[
 g=\gcd(n,3),\qquad d_0=n/g,\qquad
 M=2^{n-1},\qquad m=\binom{n+2}{3},
\]
and use the ordinary plethysm ray of P-F19,
\[
 \mu=(3),\qquad \lambda=(d_0),\qquad
 \pi=((3/g)^n).
\]
Write
\[
 f_n(s)=\dim\operatorname{Hom}_{GL_n}
 \left((\det)^{3s/g},\operatorname{Sym}^{d_0s}
                    (\operatorname{Sym}^3\mathbb C^n)\right),
 \qquad
 H_n(t)=\sum_{s\ge0}f_n(s)t^s.
 \tag{P-F31a}
\]
Let $T^{(n)}$ be its canonical Cartan multiplicity algebra,
$X=\operatorname{Spec}T^{(n)}$, and let $Y\subset X$ be the
reduced closure of the chain grading orbit. Then:

\begin{enumerate}
\item The support of the fixed scheme $X^{\boldsymbol\mu_M}$
is exactly $Y$. Away from the graded origin this fixed scheme
is the smooth chain orbit, and $X$ is smooth there. Possible
nilpotent or embedded structure at the origin is not excluded.
\item The ordinary Hilbert series $H_n(t)$ has a simple pole
at every primitive $M$th root of unity $\zeta$. More precisely,
with $g_\zeta$ the grading action,
\[
 \lim_{t\to\zeta}(1-t^M)H_n(t)
   =\frac{1}{\det(1-g_\zeta\mid N^*_{Y/X,\eta_Y})}
   \ne0.
 \tag{P-F31b}
\]
\item In every expression
\[
 H_n(t)=\frac{Q(t)}{\prod_{j=1}^{r}(1-t^{a_j})^{b_j}},
 \qquad Q(t)\in\mathbb C[t],\quad a_j,b_j\in\mathbb Z_{>0},
 \tag{P-F31c}
\]
some denominator weight $a_j$ is divisible by $M$. Thus every
such expression, in particular every positive rational form,
has maximum denominator weight at least $2^{n-1}$.
\end{enumerate}
Since the source parameter is
\[
 h=\dim GL_m+\operatorname{ht}(\lambda)+\operatorname{ht}(\pi)
   =\binom{n+2}{3}^{\!2}+1+n=O(n^6),
 \tag{P-F31d}
\]
the first, positive-form existence clause of GCT VI,
Hypothesis~1.6.7 is false for this ordinary irreducible-inner
plethysm family. The stronger canonical parameter clause was
already refuted by P-F19. No superpolynomial positivity or
saturation index is asserted here.
\end{theorem}

\begin{proof}
By P-F19 the actual canonical algebra, including multiplication,
is the entire cubic invariant ring
\[
 T^{(n)}\simeq
   \bigoplus_{s\ge0}
      \mathbb C[(\operatorname{Sym}^3\mathbb C^n)^*]^{SL_n}_{d_0s},
 \tag{P-F31e}
\]
with the coefficient grading divided by $d_0$. Scalar central
characters force every nonzero invariant degree to be a multiple
of $d_0$, so this is a regrading of the whole ring, rather than
a smaller subring. In particular $T^{(n)}$ is a finitely generated
positively graded complex domain with degree-zero part $\mathbb C$.
Its Hilbert series is exactly (P-F31a).

Consider a nonzero closed point $p$ of the fixed scheme.
P-F26 gives a representative cubic $F$ with closed $SL_n$ orbit
and an exact-form stabilizer element $A$ of finite two-power order
whose determinant has order $M$. Since $p$ is not the graded
origin, some positive-degree invariant is nonzero at $F$.
The supported cubic exponent vectors therefore span
$\mathbf1$ rationally, by the coefficient-monomial argument
in P-F27. The complete support theorem P-F30 applies to this
$F$ and $A$, without assuming that the lift has order $M$
or that the support is a graph. Its equality case identifies
$F$, after a permutation and diagonal rescaling in an eigenbasis,
with the smooth chain cubic $F_n$.

Consequently $p$ lies on the $GL_n$ orbit of the chain quotient
point. On the quotient this is exactly the grading orbit, by
(P-F26b) and the surjectivity of the character
$z\mapsto z^{-3/g}$ of $\mathbb G_m$. Conversely P-F24 shows
that the entire nonzero chain grading orbit is fixed by
$\boldsymbol\mu_M$.

The graded origin is also fixed and is in the closure of this
orbit. Indeed the map $u\mapsto u\cdot p_n$ extends from
$\mathbb G_m$ to $\mathbb A^1$: each positive-degree homogeneous
function restricts to a constant times a positive power of $u$,
and thus takes value zero at $u=0$. It follows that the closed
points of the fixed scheme are exactly those of $Y$. A
finite-type closed subset over $\mathbb C$ is determined by
its closed points, so the support of $X^{\boldsymbol\mu_M}$
is $Y$. This is an irreducible curve.

By P-F24 the ambient quotient $X$ and the fixed scheme are
smooth at every point of the nonzero chain orbit, and the fixed
scheme there has dimension one. Since these are all its
nonzero closed points, any nonreduced structure or embedded
associated point is confined to the origin. It is unnecessary
to assert reducedness at that origin: the fixed scheme has a
unique maximal-dimensional component $Y$, and $X$ is smooth
at its generic point. These are precisely the hypotheses
needed in P-F23.

Apply P-F23 with dimension $\delta=1$. It gives
\[
 \lim_{t\to\zeta}(1-t^M)H_n(t)
   =\frac{e_Y}{\det(1-g_\zeta\mid N^*_{Y/X,\eta_Y})}.
 \tag{P-F31f}
\]
All conormal eigenvalues are nontrivial. The exact grading
stabilizer and semigroup calculation in P-F28 give $e_Y=1$.
Thus the right hand side is nonzero, proving (P-F31b).
The factor $1-t^M$ has a simple zero at $\zeta$, so this is
an actual pole of order one of $H_n$.

This application includes the possible structure at the
origin. In the Koszul proof of P-F23, terms supported at
that zero-dimensional point have finite-dimensional graded
homology and hence polynomial Hilbert series. They vanish
after multiplication by $1-t^M$ in the limit. They cannot
cancel the one-dimensional generic contribution. No global
smoothness or reducedness of the fixed scheme was used.

For additional precision, P-F28 now computes the actual
normalized leading coefficient in (P-F31b). If the cyclotomic
valuation above two is normalized by $v(1-\zeta)=1$, it has
valuation
\[
 \binom{n+2}{2}-3\,2^{n-1}.
 \tag{P-F31g}
\]
The ordinary residue of $H_n(t)\,dt$ at $t=\zeta$ is
$-\zeta/M$ times the right hand side of (P-F31b), because
$(1-t^M)'|_{t=\zeta}=-M\zeta^{-1}$. The nonvanishing
in (P-F31b), rather than any comparison of phases, is
what is needed below.

Finally suppose (P-F31c) is any rational expression of
the indicated form. Its numerator is regular at $\zeta$.
If none of its denominator factors vanished there, the
expression would also be regular there, contradicting
the proved pole. Hence $\zeta^{a_j}=1$ for some $j$.
Primitivity of $\zeta$ is exactly the assertion that
$M\mid a_j$. This argument does not require a nonnegative
numerator, a prescribed number of denominator factors,
or a presentation arising from an h.s.o.p.

The exact source definition of a positive form is
GCT VI, Definition~1.2.7, printed p.~15, equation~(1.4);
its modular index is the maximum of the weights $a_j$.
Hypothesis~1.6.7, printed p.~25, first asks for the
existence of such a form with modular index polynomial
in $\dim G+\operatorname{ht}(\lambda)+\operatorname{ht}(\pi)$,
and then makes the stronger assertion for its canonical
minimal parameter form. The parameter in this family is
(P-F31d), whereas every allowed form has a weight at
least $2^{n-1}$. For any fixed polynomial in $h$, its
value at (P-F31d) is bounded by a fixed polynomial in
$n$, which is eventually smaller than $2^{n-1}$.
Therefore even the first existence clause fails.
\end{proof}

\paragraph{Consequences and remaining scope.}
The result closes the ordinary-pole gate that P-F19 alone
left open. It uses the full cubic support theorem and the
unique generically smooth fixed component; it does not
infer an ordinary pole merely from an orbit degree period,
a large parameter degree, a universal denominator, or
cohomological purity. The cancellation examples P-F21--P-F22
remain valid and explain why those shorter arguments
would not have sufficed.

Every quasipolynomial period of $f_n$ is divisible by $M$.
Indeed a quasipolynomial of period $L$ has a generating
function with denominator a power of $1-t^L$: sum its
constituent polynomials over each arithmetic progression,
using the elementary generating functions for powers.
The pole in (P-F31b) forces $M\mid L$. This does not say
that $M$ is the entire minimal period or classify its
other prime-power factors.

The implication from bounded positive-form weights to
PH2 in P-F12 remains correct, but failure of its premise
does not refute PH2 or SH. Nor do large denominator
weights alone establish a counting-complexity lower bound,
an obstruction search algorithm, or a failure of every
possible revised geometric formulation. Here the inner
shape is the single irreducible cubic shape $(3)$, and
only the number of variables grows; reducible torus
embeddings and determinant twists are not responsible
for the obstruction. No historical priority claim is
made for the support classification or the resulting
Hilbert-pole consequence.

The exact primary source for the conjecture and its
positive-form convention is
\url{https://arxiv.org/pdf/0704.0229},
Definition~1.2.7, p.~15, and Hypothesis~1.6.7, p.~25.
